\section{Introduction}
\label{sec:intro}

The mathematical capabilities of AI systems have grown at a remarkable pace.
Within a few years, they have progressed from solving grade-school word problems to matching the best human performers on competition and advanced mathematics benchmarks~\cite{hubert2026alphaproof, imo-gold, glazer2024frontiermath}.
Mathematics research, however, is a \emph{long-horizon} task.
Its goals are often not clearly defined and evolve over time. 
Success is judged not only by whether a statement is proved, but also by whether the right statement was chosen in the first place. 
It requires building a body of knowledge whose significance may be apparent only later; that is, in research, rewards are sparse, and often delayed.
% It comprises of goals that are not clearly defined, a growing body of knowledge, and sparse, often delayed reward signals. 
% Success is judged not only by whether a statement gets proved, but by whether the right statement was chosen in the first place. 
It is still an open question how to use AI effectively at this horizon.

Existing systems occupy well-specified corners of this space.
Evaluator-guided searches such as FunSearch and AlphaEvolve explore spaces of programs or constructions, but require a fixed objective and an automatic scoring function~\cite{funsearch,novikov2025alphaevolvecodingagentscientific}.
Formal-proof systems such as AlphaProof obtain a similarly clear correctness signal from a proof assistant, but require a fixed formal statement~\cite{hubert2026alphaproof}.
Interactive use of frontier models provides broad technical assistance, but leaves agenda-setting, memory, and verification with the human expert~\cite{bubeck2025science}.

In this work, we provide a detailed case study of our experience using AI operating at the level of a research \emph{program}: deciding what to attempt next, retaining what failed and why, and converting an accumulation of failures into new mathematics.
We engineered an AI research system asynchronously steered by human operators.
We directed it at a concrete open problem: the Grothendieck constant $\KG$, which quantifies the hardness threshold between hard combinatorial optimization problems and their efficient continuous relaxations, and whose value has been open since 1953.

The collaboration tightened the best known bounds on both sides,
\[
    1.7135\ldots
    \;=\;
    \frac{6\pi}{11}
    \;\le\;
    \KG
    \;\le\;
    \frac{\pi}{2\log(1+\sqrt2)} - 3.47\times10^{-4}
    \;=\;
    1.7818\ldots,
\]
The upper bound comes from a new asymptotic framework for rounding algorithms, and the lower bound from the first argument that bounds $\KG$ from below without constructing a hard instance.
Central steps of this mathematics originated with the AI system and were judged novel by domain experts; all results stated as theorems have been independently verified by the authors, and complete proofs appear in the companion paper~\cite{saha2026new}.
Our contributions are as follows.
\begin{enumerate}[label=(\arabic*), leftmargin=*]
    \item \textbf{Novel approaches for a longstanding open problem.}
    We give an expository account of the Grothendieck constant and of the results in \cref{sec:grothendieck} and \cref{sec:companion-results}.
    The upper bound introduces \emph{limiting Krivine schemes}, which enlarge the space of rounding algorithms studied since Krivine's work.
    The lower bound converts obstructions on all rounding schemes into lower bounds on $\KG$; it is the first lower bound on the constant that does not proceed by constructing a hard instance.

    \item \textbf{A methodology for long-horizon AI research.}
    We document the system in \cref{sec:ai-methodology}: its division of labor between a reasoning model and a coding agent, its file-based memory, its calibrated internal verification protocol, and its channel for asynchronous human steering.
    The run and its outcomes are quantified in \cref{sec:ai-results}.

    \item \textbf{A case study with an analysis of capabilities.}
    We reconstruct how the lower bound was found in \cref{sec:ai-case-study} and analyze the full record in \cref{sec:ai-discussion}.
    The record shows a consistent asymmetry: the system was strong at technical execution, but substantially less reliable at research judgement and at maintaining an accurate research state.
    % This suggests where human oversight has the greatest leverage, and \cref{sec:ai-future} proposes concrete directions for closing the gap.
\end{enumerate}
